\documentclass[aspectratio=169]{beamer}

% ---------------- PACKAGES ----------------
\usepackage[utf8]{inputenc}
\usepackage{graphicx}
\usepackage{booktabs}
\usepackage{amsmath}
\usepackage{xcolor}

% ---------------- THEME ----------------
\usetheme{Madrid}
\usecolortheme{default}

% ---------------- TITLE INFORMATION ----------------
\title[Smart Water Tank Quality Monitor]
{SMART WATER TANK QUALITY MONITOR}

\subtitle{IoT and Embedded Systems Project}

\author{
V. Pranathi -- 2520040104\\
B. Sahithi -- 2520040111
}

\institute{
Department of ECE
}

\date{}

% ---------------- DOCUMENT ----------------
\begin{document}

% =====================================================
% TITLE SLIDE
% =====================================================

\begin{frame}
    \titlepage
\end{frame}

% =====================================================
% ABSTRACT
% =====================================================

\begin{frame}{Abstract}

\begin{itemize}

    \item Water tanks are often monitored manually, which can lead to
    overflow, low-level conditions, and inconvenience.

    \item The proposed system uses an embedded controller and sensors
    to continuously monitor water level and temperature.

    \item Sensor data is transmitted through IoT connectivity and
    displayed on a mobile dashboard.

    \item The system enables remote monitoring and can be extended
    with alerts and additional water-quality sensors.

\end{itemize}

\end{frame}

% =====================================================
% INTRODUCTION
% =====================================================

\begin{frame}{Introduction}

\begin{itemize}

    \item Traditional water-tank monitoring usually requires
    physical inspection.

    \item Manual monitoring cannot provide continuous real-time
    information.

    \item Overflow can result in unnecessary water wastage.

    \item Temperature monitoring provides additional information
    about stored water.

    \item IoT technology allows sensor readings to be viewed
    remotely through a mobile dashboard.

    \item The project combines sensors, an embedded controller,
    wireless communication, and a user interface.

\end{itemize}

\end{frame}

% =====================================================
% PROBLEM STATEMENT
% =====================================================

\begin{frame}{Problem Statement}

\begin{itemize}

    \item Manual checking of tank level is inconvenient.

    \item Tank overflow may go unnoticed and cause water wastage.

    \item Users cannot easily monitor the tank when they are away.

    \item Water temperature is not continuously tracked.

    \item There is no simple centralized mobile view of sensor readings.

    \item \textbf{Goal:} Develop a low-cost IoT system for
    real-time tank monitoring.

\end{itemize}

\end{frame}

% =====================================================
% OBJECTIVES
% =====================================================

\begin{frame}{Objectives}

\begin{itemize}

    \item Monitor water level in real time.

    \item Measure water temperature.

    \item Process sensor readings using an embedded microcontroller.

    \item Transmit sensor data using Wi-Fi / IoT communication.

    \item Display readings on a mobile dashboard.

    \item Provide status information and support timely monitoring.

    \item Create a base for future water-quality and automatic
    pump-control features.

\end{itemize}

\end{frame}

% =====================================================
% PROPOSED SYSTEM
% =====================================================

\begin{frame}{Proposed System}

\centering

\begin{block}{System Flow}
\centering
Water Level Sensor
$\rightarrow$
ESP32 / Microcontroller
$\rightarrow$
Wi-Fi / IoT
$\rightarrow$
Cloud / Server
$\rightarrow$
Mobile Dashboard
\end{block}

\vspace{0.4cm}

\begin{itemize}

    \item \textbf{Water Level Sensor:} Measures the amount of water
    in the tank.

    \item \textbf{Temperature Sensor:} Measures stored water
    temperature.

    \item \textbf{ESP32:} Receives and processes sensor readings.

    \item \textbf{Wi-Fi / IoT:} Transfers data to the server or
    cloud platform.

    \item \textbf{Mobile Dashboard:} Displays current readings.

\end{itemize}

\end{frame}

% =====================================================
% HARDWARE AND SOFTWARE
% =====================================================

\begin{frame}{Hardware \& Software Requirements}

\begin{columns}

% LEFT COLUMN
\column{0.48\textwidth}

\textbf{Hardware}

\begin{itemize}
    \item ESP32 development board
    \item Water-level sensor
    \item Temperature sensor such as DS18B20
    \item Water tank / prototype tank
    \item Breadboard
    \item Jumper wires
    \item Power supply
\end{itemize}

% RIGHT COLUMN
\column{0.48\textwidth}

\textbf{Software}

\begin{itemize}
    \item Arduino IDE
    \item Embedded C/C++
    \item IoT platform
    \item Cloud database
    \item Mobile dashboard
\end{itemize}

\end{columns}

\end{frame}

% =====================================================
% WORKING PRINCIPLE
% =====================================================

\begin{frame}{Working Principle}

\begin{enumerate}

    \item Water-level and temperature sensors collect the tank
    readings.

    \item The ESP32 receives and processes the sensor data.

    \item The controller connects to the internet through Wi-Fi.

    \item Sensor readings are sent to the IoT platform or server.

    \item The mobile dashboard displays current water level and
    temperature.

    \item Thresholds can be configured for high- and low-level
    notifications.

\end{enumerate}

\end{frame}

% =====================================================
% EXPECTED OUTPUT
% =====================================================

\begin{frame}{Expected Output}

\begin{center}

\begin{block}{Sample Mobile Dashboard}

\Large
\textbf{WATER LEVEL}\\
78\%

\vspace{0.3cm}

\textbf{TEMPERATURE}\\
27.4$^\circ$C

\vspace{0.3cm}

\textbf{STATUS}\\
NORMAL

\end{block}

\end{center}

\end{frame}

% =====================================================
% APPLICATIONS
% =====================================================

\begin{frame}{Applications}

\begin{itemize}

    \item Smart homes and residential buildings

    \item Apartments and hostels

    \item Schools and colleges

    \item Hotels and hospitals

    \item Industrial water-storage systems

    \item Agricultural and institutional water tanks

\end{itemize}

\end{frame}

% =====================================================
% FUTURE SCOPE
% =====================================================

\begin{frame}{Future Scope}

\begin{itemize}

    \item Add a pH sensor for acidity / alkalinity monitoring.

    \item Add a turbidity sensor for water clarity.

    \item Add TDS sensing for dissolved-solids estimation.

    \item Control the water pump automatically based on level.

    \item Send mobile notifications for abnormal conditions.

    \item Store historical data and display graphs.

    \item Use data analytics or ML for water-usage prediction.

    \item Support monitoring of multiple tanks.

\end{itemize}

\end{frame}

% =====================================================
% CONCLUSION
% =====================================================

\begin{frame}{Conclusion}

\begin{itemize}

    \item The Smart Water Tank Quality Monitor provides an
    IoT-based approach to real-time tank monitoring.

    \item It combines water-level and temperature sensing with
    an embedded controller and wireless communication.

    \item A mobile dashboard enables convenient remote access
    to current readings.

    \item The system can be extended with water-quality sensors,
    alerts, and automatic pump control.

\end{itemize}

\end{frame}

% =====================================================
% THANK YOU
% =====================================================

\begin{frame}
    \centering

    \Huge
    \textbf{Thank You}

    \vspace{0.5cm}

    \Large
    Questions?

\end{frame}

% =====================================================
\end{document}